\section{Discussion}
A signed threshold can be read from a measured kernel when the allowed momentum sector, charge normalization, spin structure and metric are known. The theorem on Sol replaces the full first-gap equality problem by two integer minima for the reference metric and invariant spin lift. It does not reduce every higher signed gap to the two ground levels: period \((2,3)\) exhibits an excited-state competitor. The eta invariant is a global regularized quantity that can vanish even when these first levels split.

For a state observable one must also specify a mass, a preparation and a detector. The Nil apparatus compares two different propagators between the same projectors, so a difference is attributable to the charge within the prescribed model. A change of apparatus can change the ordering of occupations. For example at \(|u|=1/8,\ h=0,\ m_D=1\) the same lowest LC device gives approximately \(0.00823954\) and \(0.00824377\) in \(+u,-u\) order. The result supplies no universal choice of vacuum or probability law for the gravitational charge.

Several questions survive these examples. Beyond trace 60, is there a chiral cyclic monodromy with equal signed minima? For other Sol spin structures, a translated dual lattice can alter the arithmetic domain. On higher-genus Seifert bases, the missing input is the relevant twisted base Dirac spectrum. More generally a physical protocol would have to fix the charge magnitude or calibrate it independently of the sign. These questions concern different extensions of the stated sign tests.
